%----------------------------------------------------------------------------------------------------
% START OF CHAPTER
%----------------------------------------------------------------------------------------------------
\chapter{May 11}
\paragraph{Topics}
\textit{Pronouns recap, Sentence review, Conversation, Vocabulary}

\section{Grammar} \label{0511-Grammar}
\subsection{Disjunctive pronouns}
These pronouns are used as objects for prepositions like `for' \iti{per}, `between' \iti{tra},
`to' \iti{a}, and `of' \iti{di}. They are also used for emphasis and for some comparisons.
\begin{mdframed}
  \textbf{Disjunctive pronouns} pair with prepositions, not verbs. You might think of them as
  adding color or clarification to the action in a sentence. Their name describes how they
  are separated from (disjunct from) the verb.
\end{mdframed}

\paragraph{When to use disjunctive pronouns}
\begin{itemize}
    \item{With prepositions: \iti{con, da, per, senza, su, tra, ...}}
    \item{For comparison with \iti{come} and \iti{quanto} (`as/like´, `as much as')}
    \item{For emphasis}
\end{itemize}

\begin{table}[ht]
\centering
    \begin{tabular}{ r | l }
                                & pronoun    \\
        \hline
           \textbf{me}          & \iti{me}   \\
           \textbf{you}         & \iti{te}   \\
           \textbf{him}         & \iti{lui}  \\
           \textbf{her}         & \iti{lei}  \\
           \textbf{him/herself} & \iti{sè}   \\
           \textbf{us}          & \iti{noi}  \\
           \textbf{you pl.}     & \iti{voi}  \\
           \textbf{them}        & \iti{loro} \\
           \textbf{themselves}  & \iti{sè}
    \end{tabular}
    \caption{Disjunctive pronouns}
    \label{tab:tDisjunctivePronouns}
\end{table}

\paragraph{Examples}
\begin{itemize}
  \item{Are you coming to the cinema with me?} \\
  \iti{Vieni \uline{con me} al cinema?}

  \item{I talked about you to Marco.} \\
  \iti{Ho parlato \uline{di te} a Marco.}

  \item{I'm doing it for them, not for you.} \\
  \iti{Lo faccio per loro, non \uline{per te}.}

  \item{She works harder than us, it's incredible.}
  \iti{Lavora più \uline{di noi}, è incredibile.}

  \item{She's as smart as he is.} \\
  \iti{È intelligente \uline{quanto lui}.}

  \item{He called me, not you.} \\
  \iti{Ha chiamato \uline{me}, non \uline{te}.}

  \item{I live near them.} \\
  \iti{Abito vicino \uline{a loro}.}

  \item{She always talks about him.} \\
  \iti{Parla sempre \uline{di lui}.}
\end{itemize}

\subsection{Pronouns recap}

\begin{itemize}
  \item \textbf{Subject pronouns}: I, you, we, you (pl.), them $\rightarrow$ \iti{io, tu, noi, voi, loro}
  \item \textbf{Direct object pronouns}: What/whom (me, you, him/it, her/it, us, you (pl.), them) $\rightarrow$ \iti{mi, ti, lo, la, ci, vi, li, le}
  \item \textbf{Indirect object pronouns}: Recipient: to/for me, you, him, her, us, you (pl.), them $\rightarrow$ \iti{mi, ti, gli, le, ci, vi, gli}
  \item \textbf{Disjunctive pronouns}: me, you, him, her, us, you (pl.), them $\rightarrow$ \iti{me, te, lui, lei, noi, voi, loro}
\end{itemize}

\paragraph{Examples}
\begin{itemize}
  \item {I'll pay (it)} \\
  \iti{Lo pago.}

  \item {I bought him a gelato.} \\
  \iti{Gli ho comprato un gelato.}

  \item{I want to go to the store with you.} \\
  \iti{Voglio andare al negozio con te.}

  \item{Claudio went with us to the store.} \\
  \iti{Claudio è andato con noi al negozio.}
\end{itemize}

\section{Sentence review} \label{0511-Sentences}
\begin{itemize}
  \item {I told my wife to buy pizza for dinner. She told me, ``I bought it."} \\
  \iti{Ho detto mia moglie comprare la pizza per cena. Mi ha detto, ``L'ho comprata."}

  \item {My son wants ice cream this eveing. I bought it.} \\
  \iti{Mio figlio vuole il gelato questa sera. L'ho comprato.}

  \item {I didn't go to the store to buy milk and bread.} \\
  \iti{Non sono andato al negozio per comprare il latte e il pane.}

  \item {I would like to make pizza this weekend.} \\
  \iti{Vorrei fare la pizza questo fine settimana.}

  \item{I worked in my garden yesterday.} \\
  \iti{Ho lavorato nel mio giardino ieri.}

  \item{I walked with my dog today.} \\
  \iti{Ho camminato con il mio cane oggi.}

  \item{I took my dog for a walk.} \\
  \iti{Ho porto a spasso il mio cane.}

  \item{I'm buying milk because I need to have breakfast tomorrow morning.} \\
  \iti{Sto comprando latte perché ho bisogno di fare colazione domani mattina.}

  \item{I went to the store to buy milk for my wife.} \\
  \iti{Sono andato al negozio per comprare il latte per mia moglie.}

  \item{Cake? I'm sorry, I haven´t bought it because I still have gelato in the fridge.} \\
  \iti{La torta? Mi dispiace, non l’ho comprata perché io ho ancora gelato nel frigorifero.}

  \item{I haven´t bought it yet.} \\
  \iti{Non l’ho mangiata ancora.}
\end{itemize}

\section{Conversation} \label{0511-Conversation}
\begin{itemize}
  \item [F] \iti{Come state?}
  \item [A] \iti{Non c'è male.}
  \item [P] \iti{Sto bene ma, sono piuttosto stanca in questi giorni perché ho lavorato lo scorso fine settimana.}
  \item [S] \iti{Sono contento perché finalmente ho comprato la Lamborghini.}
  \item [F] \iti{Congratulazioni. Dove l’hai comprata?}
  \item [S] \iti{L’ho comprata da un amico. Un buon affare.}
  \item [P] \iti{La vorrei vedere.}
  \item [S] \iti{Potete venire con me al ristorante questa sera. Se potete.}
  \item [F] \iti{Mi dispiace. Non posso... sto ancora lavorando.}
  \item [P] \iti{Sono disponibile. A che ora?}
  \item [A] \iti{Anche io posso venire. Quale ristorante.}
  \item [S] \iti{Non ti preoccupare. Vi prendo.}
  \item [F] \iti{Che peccato! Non posso venire.}
  \item [P] \iti{La prossima volta, F.}
  \item [S] \iti{Alle otto va bene per voi?}
  \item [A] \iti{Facciamo alle otto e trenta per favore.}
  \item [P] \iti{Non è un problema.}
  \item [F] \iti{Buon divertimento!}
\end{itemize}

\section{Vocabulary} \label{0511-Vocabulary}
\subsection{Phrases}
Table \ref{tab:tPhrases-0511} has the basics.
\begin{table}[ht] %[ht] % place table roughly [h]ere or else at the [t]op of page
\centering
    \begin{tabular}{ r | l  r | l }
          I can't complain & \iti{Non c'è male}       & rather           & \iti{piuttosto}      \\
          last             & \iti{scorso}             & A good deal      & \iti{Un buon affare} \\
          available        & \iti{disponibile}        & I'll pick you up & \iti{Io vi prendo}   \\
          What a shame     & \iti{Che peccato}        & 8:30             & \iti{otto e mezza, otto e trenta} \\
          Have fun!        & \iti{Buon divertimento!} &                  &
    \end{tabular}
    \caption{Phrases}
    \label{tab:tPhrases-0511}
\end{table}

%----------------------------------------------------------------------------------------------------
% END OF CHAPTER
%----------------------------------------------------------------------------------------------------
